\section{Central-strip sources and the hierarchy}\label{app:f4-sources}
Time \(t\) counts rounds; it is an integer and is never quotiented. Fix a physical quantum
line at \((x,y,z)=(x_0,y_0,z_0)\), running along \(\sigma\), with
fold maps \(a_r\) from Lemma~\ref{f4:roles}. A cap bin has height
\(W\), bounded uniformly in the chosen window by a computable \(w\).
Its actual source hub is
\[
 \lambda=(a=x-t,\ b=y,\ \chi=\sigma+58t,\ \zeta=z+58t,
                         \ j=\lfloor t/W\rfloor),
\]
with the first four coordinates modulo \(L\) and \(j\) retained as
an integer. A hub event \(X_\lambda\) is the union of raw, intrinsic
and overlay events, with the distinct-actual-name intersection law
\(\gamma_C^{|F|}\). Suppose each wrong use has a connected lifted
witness with \(E\le c_Cg\) edges and \(g\) distinct lifted marked
hubs, bounded degree and fixed increments in all five coordinates.
Every retained bin is between zero and
\(\lfloor(H+C_{\rm stage})/W\rfloor+C_{\rm bin}\).
These inputs are verified for the actual circuit in
Section~\ref{f4:actual-sources}.

\begin{lemma}[Central-strip shadow]\label{f4:strip}
For a sufficiently large fixed increment bound \(R\), put
\(q_{\rm cut}=\lfloor L/[32R(w+1)]\rfloor\), decreasing it by one
fixed factor if necessary to include root-use spurs. At all sufficiently
large \(L/(w+1)\), one source shadow agrees with every quantum root
and has fixed Lipschitz increments on every rooted witness with fewer
than \(q_{\rm cut}\) edges.
\end{lemma}
\begin{proof}
Choose the integer \(\tau(\lambda)\equiv x_0-a\pmod L\) nearest to
\((j+\tfrac12)W\), with a fixed tie rule, and set
\[
 \pi_r(\lambda)=
   (a_r(\chi-58\tau(\lambda)),\tau(\lambda)).
\]
At an exact quantum root \((\sigma,t)\), the time \(t\) has that
congruence and belongs to its bin. Since \(W<L\), it is the unique
nearest representative. The cut in round \(t\) uses the first coordinate computed in that
round. If that output's hub is indexed by the next boundary, join it
to the quantum-use root by its fixed one-round spur before counting
edges. This is the timing convention throughout the proof.

Lift a spanning tree from a root and put
\(t_v=t_{\rm root}-(a_v-a_{\rm root})\). With increment bounds
\(R_a,R_j\le R\), a path of \(i<q_{\rm cut}\) edges obeys
\[
 |t_v-(j_v+\tfrac12)W|
 \le W/2+i(R_a+WR_j)
 \le w/2+q_{\rm cut}R(w+1)<L/8.
\]
A bounded root spur is absorbed by the fixed reduction of the cutoff
and a lower-size threshold. Thus this lift gives the unique midpoint
representative at every vertex. Two paths cannot have nonzero
\(L\)-monodromy in \(a\), since their total displacement is below
\(2Rq_{\rm cut}<L\). The same reasoning identifies repeated lifted
occurrences projecting to one actual hub. On each selected edge,
\(|\Delta\tau|\le R_a\), and the physical \(\sigma\)-shadow changes
by at most \(R_x+58R_a\). Lemma~\ref{f4:roles} preserves this bound.
The globally defined midpoint map is therefore consistent on all
selected short witnesses, with the same value for each actual name.
Only those edges are used; edges crossing an arbitrary representative
branch cut elsewhere in the quotient graph need not be Lipschitz.
\end{proof}

\begin{lemma}[Long cap witnesses]\label{f4:cap-exception}
For sufficiently small fixed \(\gamma_C\), the event that any needed
cap explanation through \(H\) has at least \(q_{\rm cut}\) edges
has probability at most
\[
 C_{\rm cap}(H+C_{\rm cap})L^4
           \exp[-\beta_{\rm cap} L/(w+1)].
\]
Outside this event every needed cap use has a complete short certificate
with distinct actual names and the shadow of Lemma~\ref{f4:strip}.
\end{lemma}
\begin{proof}
Increase \(c_C\ge1\) for isolated roots and fixed stars. A connected
\(E\)-edge witness has a spanning-tree depth-first walk of length at
most \(2E\). Assign each distinct lifted mark to its first visit and
partition the walk into pieces of at most \(q_{\rm cut}\) edges,
retaining shared boundary vertices. For \(E\ge q_{\rm cut}\) the
number of pieces is at most \(2E/q_{\rm cut}+1\le3c_Cg/q_{\rm cut}\).
Some piece therefore contains at least \(q_{\rm cut}/(3c_C)\)
assigned distinct lifted marks. For a sufficiently reduced fixed cutoff
its characteristic diameters are below \(L/2\). Distinct lifts of
one actual hub have the same integer bin and differ by an \(L\)-deck
vector in a spatial characteristic. They cannot both occur in this
piece. Its marks are distinct actual names before their law is applied.

The causal bin bound supplies at most \(C(H+C)L^4\) starting vertices.
Finite degree, roles and mark subsets give at most
\(K_{\rm cap}^{q_{\rm cut}+1}\) descriptions per start. For each
admissible description apply the actual source law to its distinct
marks; other codes only overcount. Hence the exceptional probability
is at most
\[
 C(H+C)L^4 K_{\rm cap}^{q_{\rm cut}+1}
             \gamma_C^{\lfloor q_{\rm cut}/(3c_C)\rfloor}.
\]
Choose the source cap after \(c_C,K_{\rm cap},R\) to make its
exponential base less than one. Floors and small cutoffs are absorbed
in a fixed threshold and prefactor. This proves the bound. Every
remaining complete explanation is short, and the same diameter
argument makes all its marked names injective. Select a minimal one
if several are available. The estimate is unconditional.
\end{proof}

\begin{lemma}[Mixed roles and gluing]\label{f4:mixed}
Let raw hardware/calendar names have joint law \(p_H^{|F|}\), with
one actual hub for every use of the same site and round. Duplicate
cap and raw names into two formal fold roles, retaining their underlying
events. For \(g_C,g_H\) distinct formal names of the two families,
\[
 \Pr(\hbox{all requested events})
 \le\gamma_C^{g_C/4}p_H^{g_H/4}
 \le\gamma_{\rm eff}^{g_C+g_H},\qquad
 \gamma_{\rm eff}=\max(\gamma_C^{1/4},p_H^{1/4}).
\]
Gluing complete cap and raw witnesses at their actual quantum-use
roots and adding the dispatcher stars gives bounded degree, fixed
edge/source size and one Lipschitz shadow on all admitted short
certificates.
\end{lemma}
\begin{proof}
Each family names at least half as many actual hubs as formal names.
Apply its joint law, then Cauchy--Schwarz between the two possibly
dependent families. This gives the fourth powers displayed above.
Raw calendar errors and direct hardware uses are annotations of one
raw hub; raw/statistical/overlay cap types and old/new occurrences
remain annotations of one cap hub.

For a fixed cap coordinate and bin at the quantum line, \(a=x_0-t\)
determines time modulo \(L\), and \(W<L\) permits at most one such
time in the bin. The remaining coordinates fix the quantum root.
The gluing degree is consequently bounded; fixed root-use spurs only
enlarge this bound finitely. Shadows agree at the common root by the
timing convention of Lemma~\ref{f4:strip}. Raw shadows use literal
\((\sigma,t)\) followed by \(a_r\). Dispatcher attachment stars have
fixed increments in each formal role, including the unfinished-module
station envelope. Existing source names are retained throughout.
The claimed size, entropy and shadow properties follow by enlarging
these fixed catalogue constants once. Shared physical flags across
two roles have already been paid by the fourth-root weight.
\end{proof}

\begin{lemma}[Alias-safe source recursion]\label{f4:recursion}
Suppose the scheduled circuit has complete regular boxes of length
\(P^\ell\), time height \(T^\ell\), fixed child-index halos, at most
\(A_{\rm rec}\) candidate pairs, and two full bad child supports
separated by \(\rho_{\rm gap}P^\ell\) whenever a parent is bad.
Root counts are polynomial in \(P^\ell\). Choose \(s_{\rm top}\)
larger than the complete halo and \(s_{\rm top}P^k\le L/2\).
At a sufficiently small fixed source cap, short-certified bad-box
events under the original physical law satisfy
\[
 \Pr(N_\ell^{\rm sh})\le z_\ell\le D\,2^{-2^\ell},\qquad
 D=1/(2\bar A),\quad \bar A=\max(1,A_{\rm rec}),\quad0\le\ell\le k.
\]
The exceptional events of Lemmas~\ref{f4:raw} and
\ref{f4:cap-exception} are paid separately.
\end{lemma}
\begin{proof}
A base root is short-certified if it has one of the admitted complete
certificates. Define the geometric bad-box predicate from these roots
under the original law. Outside the two exceptions this includes all
actual bad base roots. Work with sums of source weights of fully
embedded, recursively pruned certificates, as in Lemma~\ref{d:recursion}.
At a parent of level-\(\ell\) children, if a primitive witness has full
graph length
at least \(\epsilon_{\rm sh}P^\ell\), discard its entire repeated
subtree and pay that witness once. Bounded anchored entropy and the
joint source law bound this contribution by
\(B_{\rm tail}P^{d_0\ell}e^{-a_{\rm tail}P^\ell}\).
Choose \(\epsilon_{\rm sh}\) so a shorter witness has shadow reach
below one third of the separated-support gap. The selected two child
certificates then have disjoint formal source-name sets, because every
name has one shadow. Their joint weight is the product of their
weights. Within a child apply the same pruning recursively. Full halos
are local cycle arcs by the choice of \(s_{\rm top}\); time stays
unquotiented. The resulting sum satisfies
\[
 z_{\ell+1}\le A_{\rm rec}z_\ell^2+
       B_{\rm tail}P^{d_0\ell}e^{-a_{\rm tail}P^\ell}.
\]
This multiplies weights of disjoint names, not marginal bad-box
probabilities. A witness longer than the scale threshold is still
counted at its full length. If that threshold exceeds the cap cutoff,
there are simply no admitted cap witnesses in that case.

Choose the fixed cap so that \(z_0\le1/(4\bar A)\) and
\[
 a_{\rm tail}\ge\max(0,\log(4\bar A B_{\rm tail}))
                 +d_0\log P/P+2\log2.
\]
Exactly the comparison in Lemma~\ref{d:recursion} makes the remainder
at most \((D/2)2^{-2^{\ell+1}}\); the squared term has the same bound
since \(A_{\rm rec}D^2\le D/2\). Induction proves the claimed form,
with exponent \(\log2\), uniformly in the window. None of these laws
is conditioned on absence of the exceptional events.
\end{proof}
